%------------------------------------------------------------------------------%
% Luis A. D'Afonseca
%
%------------------------------------------------------------------------------%
\documentclass[a4paper,12pt,fleqn]{article}

\usepackage{style_avaliacao}
\usepackage{systeme}

\ShowAnswers

\newcommand{\D}[2]{\frac{\partial {#1}}{\partial {#2}}}
\newcommand{\DS}[2]{\frac{\partial^2 {#1}}{\partial {#2}^2}}
\newcommand{\DM}[3]{\frac{\partial^2 {#1}}{\partial {#2} \partial {#3}}}

%------------------------------------------------------------------------------%
\begin{document}

\makeHeader{IS}{Prova 3 -- Turma 8}{10/09/2024}

% 01
%------------------------------------------------------------------------------%
\exercise{25}
Determine se a série
\(
  \sum_{n=1} \frac{n-2}{n^3-n^2+3}
\)
covnerge ou diverge.
Explicite o teste usado.

\begin{answer}
  Vamos usar o teste da comparação no limite, comparando com a série
  \[
    \sum_{n=1} \frac{1}{n^2}
  \]
  que é uma série convergente pois é uma série $p$ com $p=2>1$.

  \vspace{\baselineskip}
  Calculando o limite
  \begin{align*}
    L
    & = \lim_{n\to\infty} \frac{n-2}{n^3-n^2+3} \left(\frac{1}{n^2}\right)^{-1} \\[2mm]
    & = \lim_{n\to\infty} \frac{(n-2)n^2}{n^3-n^2+3} \\[2mm]
    & = \lim_{n\to\infty} \frac{n^3-2n^2}{n^3-n^2+3} \\[2mm]
    & = \lim_{n\to\infty} \frac{1-\sfrac{2}{n}}{1-\sfrac{1}{n}+\sfrac{3}{n^3}} \\[2mm]
    & = 1
  \end{align*}

  Como o $L=1>0$ a série
  \(
    \sum_{n=1} \frac{n-2}{n^3-n^2+3}
  \)
  converge.
  \clearpage
\end{answer}

% 02
%------------------------------------------------------------------------------%
\exercise{25}
Determine o interior do intervalo onde a série
\(
  \sum_{n=1}^\infty \frac{(-1)^n(x+2)^n}{n}
\)
converge.

\begin{answer}
  Vamos aplicar o teste da razão (seria possível aplicar o teste da raiz também).

  O termo geral da série é
  \[
    u_n = \frac{(-1)^n(x+2)^n}{n}
  \]
  seu sucessor é
  \[
    u_{n+1} = \frac{(-1)^{n+1}(x+2)^{n+1}}{n+1}
  \]
  Calculando o quociente dos módulos
  \begin{align*}
    \frac{\abs{u_{n+1}}}{\abs{u_n}}
    & = \abs{\frac{(-1)^{n+1}(x+2)^{n+1}}{n+1}} \; \abs{\frac{n}{(-1)^n(x+2)^n}} \\[2mm]
    & = \frac{\abs{x+2}^{n+1}}{n+1} \; \frac{n}{\abs{x+2}^n} \\[2mm]
    & = \frac{\abs{x+2}n}{n+1}
  \end{align*}
  Calculando o limite
  \[
    \rho
    = \lim_{n\to\infty} \frac{\abs{u_{n+1}}}{\abs{u_n}}
    = \lim_{n\to\infty} \frac{\abs{x+2}n}{n+1}
    = \abs{x+2} \lim_{n\to\infty} \frac{n}{n+1}
    = \abs{x+2}
  \]
  A condição de convergência do teste da razão é $\rho<1$ portanto
  \begin{align*}
    \rho & < 1 \\
    \abs{x+2} & < 1 \\
    -1 < x+2 & < 1 \\
    -1 -2 < x & < 1 - 2 \\
    -3 < x & < -1
  \end{align*}
  A série converge absolutamente no intervalo \((-3, -1)\).
  \clearpage
\end{answer}

% 03
%------------------------------------------------------------------------------%
\exercise{25}
Calcule a integral indefinida da função
\(
  f(x) = \sum_{n=0}^\infty \frac{(n!)^2x^n}{2^n(2n)!}
\)

\begin{answer}
  \begin{align*}
    I
    & = \int \left(\sum_{n=0}^\infty \frac{(n!)^2x^n}{2^n(2n)!}\right) dx \\[2mm]
    & = \sum_{n=0}^\infty \left(\int \frac{(n!)^2x^n}{2^n(2n)!} dx\right) \\[2mm]
    & = \sum_{n=0}^\infty \frac{(n!)^2}{2^n(2n)!} \int x^n dx \\[2mm]
    & = C + \sum_{n=0}^\infty \frac{(n!)^2}{2^n(2n)!} \frac{x^{n+1}}{n+1} \\[2mm]
    & = C + \sum_{n=0}^\infty \frac{(n!)^2x^{n+1}}{2^n(n+1)(2n)!}
  \end{align*}
  \clearpage
\end{answer}

% 04
%------------------------------------------------------------------------------%
\exercise{25}
Encontre a Série de Taylor, centrada em 2, da função \( f(x) = e^x \)

\begin{answer}
  Calculando as derivadas de
  \( f(x) = e^x\)
  temos
  \[
    f^{(n)}(x) = e^x
  \]
  Avaliando em $x=2$ temos
  \[
    f^{(n)}(2) = e^2
  \]
  O coeficiente da Série de Taylos será
  \[
    c_n
    = \frac{f^{(n)}(2)}{n!}
    = \frac{e^2}{n!}
  \]
  A Série de Taylor é
  \[
    \sum_{n=0}^\infty c_n (x-1)^n
    = \sum_{n=0}^\infty \frac{e^2}{n!} (x-2)^n
  \]
\end{answer}

%------------------------------------------------------------------------------%
\end{document}
%------------------------------------------------------------------------------%
